\documentclass[11pt,twocolumn]{article}
\usepackage[utf8]{inputenc}
\usepackage{amsmath,amssymb,amsfonts}
\usepackage{geometry}
\usepackage{booktabs}
\usepackage{hyperref}
\usepackage{cite}
\usepackage{microtype}
\geometry{letterpaper, margin=1in}

\title{\textbf{The Knife and the Actuator: Language Games, Direction of Fit, and the Operative Limits of Description}}
\author{\textbf{Watson Hartsoe}\\
\small Operative Humanities Laboratory & Worldful Press\\
\small \texttt{Working Paper · AI-augmented research process}}
\date{September 7, 2026}

\begin{document}
\maketitle

\begin{abstract}
Autoregressive language models and agentic software frameworks operate under an unexamined assumption of semantic symmetry: the belief that because dynamic systems can be described in text, they can be reliably controlled through text. This paper argues that this premise collapses when confronted with the fundamental asymmetry of direction of fit (Anscombe 1957; Searle 1975). A description bears a word-to-world direction of fit ($d_\downarrow$), wherein errors demand symbolic revision. An actuation bears a world-to-word direction of fit ($d_\uparrow$), wherein errors demand physical transformation of the external environment. When an autonomous software agent emits tokens intended to actuate external environments—whether dispatching shell commands, routing cybernetic workflows, or steering mechanical actuators—the operation undergoes a violent dimensional reduction: what we formalize as the knife schema (Hartsoe 2026). Drawing upon 44 atomic research objects compiled across 12 canonical language games—spanning speech act theory (Austin, Searle), situated robotics (Agre & Chapman, Suchman), organizational commitment state machines (Winograd & Flores), indeterminate scoring (Cardew, Cage), and constraint mechanics (Oulipo, Elster)—we diagnose the ekphrastic illusion of contemporary AI: the fatal confusion of narrative coherence with mechanical compliance. We demonstrate that robust agency does not emerge from constructing monolithic world models or lengthening chain-of-thought descriptions, but from indexical-functional coupling to substrate resistance. Where description terminates, the knife begins; what survives the crossing is not an unblemished semantic map, but the physical remainder of the cut.
\end{abstract}

\noindent\textbf{Keywords:} Direction of Fit; Language Games; Ekphrasis; Situated Action; Autoregressive Agents; Deictic Representation; Substrate Resistance; Speech Acts.

\section*{1. Introduction: The Asymmetry of the Counter-Cut}
In contemporary artificial intelligence engineering, text has been quietly crowned as a universal actuator. In frameworks governing large language model agents, tool-use APIs, and chain-of-thought planners, the emission of a sequence of tokens is routinely treated as isomorphic to physical action. A prompt produces an execution string; the string is passed to an interpreter; a file is modified, a motor fires, or an API call alters a database record. Because the entire pipeline is mediated through linguistic tokens, system designers operate under the tacit assumption of semantic continuity: that the boundary between describing an action and executing it is merely a matter of syntactical formatting.

This paper contends that this assumption commits a foundational category error. The boundary between description and actuation is not a smooth gradient of increasing specificity; it is a topological rupture. Drawing upon the aperture treatises of Hartsoe (2026), we term this rupture the knife. A knife does not represent the matter it divides; it alters its physical topology. The linguistic representation of a cut operates within an unconstrained, high-dimensional combinatorial space where contradictory sentences can peacefully coexist on adjacent lines. The physical cut, by contrast, is an irreversible, dimensionally reductive intervention in continuous material reality. When the knife descends, the matter separates, but the blade simultaneously absorbs mechanical wear: the unmodelled counter-cut.

The failure to recognize the knife produces what we term the ekphrastic illusion of autonomous agency. In classical rhetoric, ekphrasis refers to the poetic representation of visual art—the substitution of vivid linguistic description for the physical presence of an object. In computational systems, the ekphrastic illusion arises when an architecture mistakes the internal syntactic coherence of an agent's self-generated narrative (its reasoning trace or world model) for causal compliance on the part of the external environment. When reality resists the model, the typical engineering response is to demand more ekphrasis: larger context windows, denser retrieval-augmented graphs, and longer chain-of-thought deliberations. Yet no quantity of description can cross the threshold of the aperture to force mechanical compliance. To understand why, we must examine the formal mechanics of direction of fit.

\section*{2. The Mechanics of Direction of Fit: Anscombe, Searle, and the Inversion of Mistake}
The formal distinction between description and actuation was established by G. E. M. Anscombe (1957) through her canonical parable of the shopping list and the detective's log. Consider a man walking through a market with a shopping list written on a slip of paper. Behind him walks a detective, recording every item the man places in his basket. At the end of the journey, both slips of paper contain identical sequences of inscriptions: "butter, milk, bread." Syntactically, the two lists are indistinguishable. Yet their operational semantics are inverted. If the man buys margarine instead of butter, the mistake lies in his performance; he cannot rectify the error by crossing out "butter" and writing "margarine" on his list. The list is an invariant directive that obligates the world to conform to the ink. If the detective writes "margarine" when the man bought butter, the mistake lies in the record; the detective cannot rectify the error by forcing the man to replace the butter. The record is an assertive report that is obligated to conform to the world.

John Searle (1975) formalized Anscombe's insight into the foundational taxonomy of illocutionary force, defining the direction of fit as a mathematical vector between the symbolic sign ($W$) and the physical world ($I$). In an assertive speech act ($d_\downarrow$, word-to-world), the speaker's illocutionary point is to match an independent reality. In a directive or commissive speech act ($d_\uparrow$, world-to-word), the illocutionary point is to alter reality so that it matches the propositional content of the token. 

This asymmetry reveals the catastrophic flaw in treating autoregressive token generators as autonomous deciders. An autoregressive language model is trained under a strictly assertive loss function: next-token cross-entropy error measures whether the emitted string matches the statistical distribution of prior inscriptions ($d_\downarrow$). When such a model is embedded into an agent loop and instructed to "plan" or "actuate," the runtime environment abruptly demands that the generated tokens function as directives ($d_\uparrow$). But the model possesses no internal mechanism to distinguish whether a discrepancy between its output and the environment constitutes an error in its performance or an error in its description.

This confusion is compounded by what J. L. Austin (1962) diagnosed as the doctrine of infelicities. Performative utterances—actions performed through words, such as promising, christening, or ordering—do not succeed or fail based on truth-conditions. They succeed or fail based on felicity conditions: the institutional, physical, and relational authority of the speaker within an accepted conventional procedure. When an agent emits the token "I promise to refund your payment," but lacks the legal status, bank credentials, or institutional sanction to transfer funds, the utterance is not a false statement; it is an infelicitous misfire ($A.1$). In the terminology of Winograd and Flores (1986), language in human activity does not function as an abstract transmission of information, but as a commitment machine—an explicit state network of requests, promises, declines, and declarations governed by the possibility of breakdown. When autonomous agents operate without explicit commitment state machines, they mistake the syntactical generation of a promise for its institutional fulfillment.

\section*{3. Ekphrasis vs. Actuation: The Illusion of Symmetric Semantic Mediation}
The ekphrastic illusion is seductive because within digital environments, the illusion of symmetry is maintained by software virtualization. In a sandboxed software container, deleting a file or updating a variable appears trivial and cleanly reversible. This leads system architects to treat all external domains as compliant computational memory. But the moment computation interfaces with external sociotechnical or physical substrates, the smooth symmetry of symbolic manipulation evaporates.

Consider musical notation and scoring, as explored by Nelson Goodman (1968), John Cage (1961), and Cornelius Cardew (1971). Goodman demonstrated that a musical score is allographic: it defines a compliance class of performances such that any acoustic performance that satisfies the specified pitch and duration requirements constitutes an authentic instance of the work. The score operates under the presumption that performers share an unambiguous, discrete notation system. Yet when Cardew composed *Treatise* (1967–1971)—a 193-page graphic score consisting entirely of abstract lines, circles, and geometric diagrams devoid of musical notation—he explicitly severed the allographic link. Cardew recognized that a graphic score does not instruct the performer what sound to produce; rather, it forces the performer to confront the physical resistance of their instrument and the social dynamics of the ensemble. The score is not an algorithm for sound, but an operational ekphrasis that provokes an unscripted material performance.

When modern AI systems construct long chain-of-thought explanations prior to dispatching an API call, they are generating Cardew-like graphic scores while assuming they possess Goodman-like allographic compliance. The agent's chain-of-thought is a rhetorical performance designed to satisfy the user's demand for explainability. But the receiving actuator (the database driver, the industrial robot, or the human collaborator) does not execute the reasoning trace; it receives only the terminal impulse. As Hartsoe (2026) observes in *Listener with Actuators*, the listener wired to mechanical actuators is fundamentally subordinate to the physical physics of its relays. If an industrial arm is commanded to move to coordinate $(X, Y, Z)$ at velocity $V$, the physical motor knows nothing of the elaborate ethical or causal rationale that preceded the command. The cut is made; the relay trips; the dimensional collapse is instantaneous.

\section*{4. Deictic Coordination and the Rejection of Global World Models}
The standard computational response to environmental failure has historically been the construction of increasingly comprehensive world models. In classical artificial intelligence planning, exemplified by STRIPS (Fikes and Nilsson 1971), an agent maintains an exhaustive predicate-calculus model of the world and performs heuristic search over state-space graphs to prove that a sequence of operators transforms the initial state into the goal state. 

This classical planning paradigm was decisively dismantled by Philip Agre and David Chapman (1987) in their implementation of Pengi, and by Lucy Suchman (1987) in her ethnomethodological critique of situated action. Agre and Chapman demonstrated that an agent operating in a hostile, dynamic video-game environment cannot survive by constructing and searching global world models. Centralized search trees suffer from fatal combinatorial latency: by the time the agent calculates an optimal path, the moving projectile has already destroyed it. Instead, Agre and Chapman pioneered deictic representations—indexical-functional entities that couple sensory focus directly to immediate behavioral impulses. Rather than representing an objective entity such as "Bee-34 at coordinates $(14, 22)$," the agent represents "the-bee-I-am-running-away-from." Deictic representations are relative to the agent's immediate physical posture and current project, executed through combinational logic circuits without memory or search.

Suchman (1987) proved that human action in dynamic environments follows an identical situated logic. Plans are not prospective programs that control action; they are retrospective linguistic resources used to explain, justify, or reconstruct action after the fact. When human beings navigate complex, real-world tasks—such as operating a photocopier or navigating a naval vessel (Hutchins 1995)—coordination is achieved not through internalizing global symbolic maps, but through situated interactions with physical artifacts, gestures, and environmental affordances. 

As Adam Kendon (2004) and David McNeill (1992) have documented in the study of gesture, human bodily communication does not rely on a single, uniform linguistic code. Speech is linear and segmented; gesture is global and synthetic. In McNeill's formulation of the "growth point," utterance production begins not with an abstract sentence outline, but with a dialectical clash between an indexical, imagistic bodily gesture and a conventional linguistic category. To reduce all agentic intelligence to sequential textual tokens is to discard the entire indexical-spatial substrate of real-time coordination.

\section*{5. Substrate Resistance and the Incomputable Remainder}
What is systematically erased by the ekphrastic illusion is the phenomenon of substrate resistance. Every expressive medium, every mechanical apparatus, and every institutional structure possesses an inalienable material friction that resists symbolic totalization.

In the sociology of technology, Bruno Latour (1996) chronicled this resistance in *Aramis, or the Love of Technology*. Aramis—an automated, personal rapid transit system conceived in Paris—failed not because its mathematical algorithms were erroneous, but because its engineers believed the pristine elegance of their blueprint could impose itself upon human politicians, electrical couplings, and suburban municipalities without negotiation. The designers refused to let the project be contaminated by the messy compromises demanded by its non-human and human actors. Aramis remained a pure ekphrasis, and because it refused to become a knife that accepted the wear of the cut, it died in a warehouse.

A parallel lesson emerges from art history and textual scholarship. Michael Baxandall (1972) demonstrated that the Florentine painting contract of the fifteenth century was not a purely aesthetic mandate, but a commercial instrument designed to tether artistic imagination to material scarcity: the contract explicitly stipulated the precise grade of ultramarine blue and gold leaf to be used, anchoring symbolic representation in physical cost. In editorial theory, John Bryant (2002) and Jerome McGann (1983) established that no text is an immutable, ideal Platonic form; every text is a "fluid text," continually revised and deformed by the material apparatus of printing presses, scribal errors, editorial censorship, and reader intervention.

This material friction is not a flaw to be engineered away; it is the very engine of generativity. This truth was demonstrated with mathematical purity by the Oulipo (Ouvroir de Littérature Potentielle), founded by Raymond Queneau and François Le Lionnais (Motte 1986). When Georges Perec (1969) wrote *La Disparition*, a 300-page novel without a single occurrence of the letter 'e', the constraint was not a passive limitation. The systematic exclusion of the most common vowel in the French language forced Perec to abandon conventional idioms and unearth latent lexical possibilities that would have remained invisible in unconstrained prose. Furthermore, as Queneau recognized, an absolute structural constraint risks decaying into sterile autopoiesis unless it incorporates the *clinamen*—the deliberate, minute infraction of the rule that reintroduces the contingency of the physical world.

When an autonomous agent operates in a real environment, substrate resistance manifests as the counter-cut. When the blade divides the wood, the grain of the wood deflects the blade. The deflection cannot be calculated from the blueprint because the cellular structure of that specific plank of timber is unique to its biological history. The deflection is the incomputable remainder of actuation.

\section*{6. Conclusion: Toward an Operative Pragmatics of the Aperture}
The contemporary rush to deploy autonomous language agents across commercial, scientific, and industrial infrastructure has reached an impasse defined by hallucination, fragility, and prompt decay. We have argued that these pathologies are not temporary benchmark deficits to be solved by scaling model parameters or expanding token contexts. They are the inevitable symptoms of the ekphrastic illusion—the structural failure to distinguish word-to-world description from world-to-word actuation.

To move beyond this impasse, artificial intelligence must abandon the dream of all-purpose linguistic virtualization and develop an operative pragmatics of the aperture. Such a pragmatics must be grounded in four structural principles:

1. **Bifurcation of Direction of Fit:** Architectures must formally separate assertive logging registers ($d_\downarrow$) from directive actuation channels ($d_\uparrow$). A model must never be allowed to evaluate its own execution success solely through subsequent linguistic self-reflection; success can only be determined by non-linguistic sensors monitoring physical state changes.
2. **Deictic Primacy over Global Mapping:** Rather than constructing massive, persistent vector databases that simulate global world states, autonomous agents must utilize indexical-functional representations (Agre & Chapman 1987) that couple immediate perceptual markers to reflexive execution logic.
3. **Conversational Breakdown as Native State:** Following Winograd and Flores (1986), systems must treat breakdown not as an exception or runtime error, but as the foundational mode of cooperative communication. An agent must possess explicit state machines for declaring breakdowns, renegotiating commitments, and acknowledging lack of authority.
4. **Integration of the Counter-Cut:** Systems must measure and model the wear on their own instruments. When an action fails, the failure must be registered not as a prompt prompting error, but as evidence of substrate friction that demands adaptation of the physical tool.

Where description terminates, the knife begins. What survives the crossing of the aperture is never the complete, pristine semantic dream of the prompt, but the tangible, scarred, and consequential mark of the cut.

\bibliographystyle{plain}
\bibliography{20260907-slipcase__references}

\appendix
\section*{Appendix A: Assembly Instrument Pointer}
This publication was assembled under the POML v15.55-AM Lossless Research Compiler Specification. The full assembly prompt is preserved verbatim in \texttt{_PROMPTS/PROMPT\_\_POML\_15\_55\_AM\_\_SLIPCASE\_COMPILATION.txt} and embedded in \texttt{index.html}.

\section*{Appendix B: Making History Summary}
\textbf{Provided:} 44 atomic zettel cards across 12 language games.\\
\textbf{Preserved:} 100\% exact payloads, unique SHA-256 hashes, BibTeX blocks.\\
\textbf{Verified:} Zero broken links, zero duplicate payloads, local Chrome PDF compilation.\\
\textbf{Unverified:} Physical long-term archival print endurance.

\section*{Appendix C: Replication Path}
To verify, inspect, or rebuild this field, open \texttt{index.html} or run \texttt{python3 WAYS TO WRITE/slipcase_builder/01_cards_and_graph.py}. Rejoin phrase: \textit{``where description terminates, the knife begins''}.

\end{document}
